\subsection{FUNCTORIAL PROPERTIES OF THE LIE ALGEBRA}

Let \(G\) and \(H\) be Lie groups and \(\phi\) a morphism of \(G\) into \(H\) . The restriction of \(U(\phi)\) to \(L(G)\) , which is just \(T_e(\phi)\) , is a continuous morphism of \(L(G)\) into \(L(H)\) , which we denote by \(L(\phi)\) . If \(\psi\) is a morphism of \(H\) into a Lie group, then \(L(\psi \circ \phi) = L(\psi) \circ L(\phi)\) .

For \(\phi\) to be an immersion, it is necessary and sufficient that \(\mathrm{L}(\phi)\) be an isomorphism of \(\mathrm{L}(\mathrm{G})\) onto a subalgebra of \(\mathrm{L}(\mathrm{H})\) admitting a topological supplement. In particular, if \(\mathrm{G}\) is a Lie subgroup of \(\mathrm{H}\) and \(\phi\) is the canonical injection, \(\mathrm{L}(\mathrm{G})\) is identified with a Lie subalgebra of \(\mathrm{L}(\mathrm{H})\) by means of \(\mathrm{L}(\phi)\) . More particularly, if \(\mathrm{G}\) is an open subgroup of \(\mathrm{H}\) , \(\mathrm{L}(\mathrm{G}) = \mathrm{L}(\mathrm{H})\) .

If G is a Lie quasi-subgroup of H, L(G) is also identified with a closed Lie subgroup of L(H).

For \(\phi\) to be a submersion, it is necessary and sufficient that \(\mathrm{L}(\phi)\) be surjective and that its kernel admit a topological supplement. In that case, the kernel \(N\) of \(\phi\) is a Lie subgroup of \(G\) and \(\mathrm{L}(N) = \mathrm{Ker} \, L(\phi)\) . In particular, if \(H\) is the quotient Lie group of \(G\) by a normal Lie subgroup \(P\) , \(L(P)\) is an ideal of \(\mathrm{L}(G)\) and, if \(\phi\) is the canonical surjection of \(G\) onto \(H\) , \(\mathrm{L}(G/P)\) is identified with \(\mathrm{L}(G)/\mathrm{L}(P)\) by means of the morphism derived from \(\mathrm{L}(\phi)\) when passing to the quotient.

Let \(I\) be a finite set, \((\mathbf{G}_{i})_{i\in I}\) a family of Lie groups, \(G\) their product and \(p_i\) the canonical morphism of \(G\) onto \(G_i\) . Then \((\mathrm{L}(p_i))_{i\in I}\) is a morphism of the Lie algebra \(L(G)\) into the Lie algebra \(\prod_{i\in I} L(G_i)\) and is an isomorphism of normable spaces. \(L(G)\) is therefore identified with \(\prod_{i\in I} L(G_i)\) by means of \((L(p_i))_{i\in I}\) .

PROPOSITION 28. Let G and H be Lie groups and \(\phi\) a morphism of G into H. Suppose that K is of characteristic 0 and that H is finite-dimensional.

\begin{enumerate}
\def\labelenumi{(\roman{enumi})}
\item
  The kernel N of \(\phi\) is a Lie subgroup of \(\mathbf{G}\) and \(\mathrm{L}(\mathbf{N}) = \mathrm{Ker} \, \mathrm{L}(\phi)\) .
\item
  The morphism \(\psi\) of G/N into H derived from \(\phi\) when passing to the quotient is an immersion.
\item
  If \(\phi(\mathbf{G})\) is closed in \(\mathbf{H}\) and the topology of \(\mathbf{G}\) has a countable base, \(\phi(\mathbf{G})\) is a Lie subgroup of \(\mathbf{H}\) , \(\psi\) is an isomorphism of the Lie group \(\mathbf{G}/\mathbf{N}\) onto the Lie group \(\phi(\mathbf{G})\) and \(\mathbf{L}(\phi(\mathbf{G})) = \mathbf{Im} \mathbf{L}(\phi)\) .
\end{enumerate}

Let G operate on H on the left by the mapping \((g, h) \mapsto \phi(g)h\) . It suffices to apply Proposition 14 of §1, no. 7, to the orbit of e.

PROPOSITION 29. Let G and H be Lie groups and \(\phi\) a morphism of G into H. Suppose that K is of characteristic 0 and that H is finite-dimensional. If H' is a Lie subgroup of H, then G' = \(\phi^{-1}(\mathrm{H}')\) is a Lie subgroup of G and L(G') = L(φ)\(^{-1}\)(L(H')).

Let \(\pi\) be the canonical mapping of H into the homogeneous space \(\mathbf{X} = \mathbf{H} / \mathbf{H}'\) . Let G operate on \(\mathbf{X}\) on the left by the mapping \((g, x) \mapsto \phi(g)x\) . The stabilizer of \(\pi(e)\) is \(G'\) , which is therefore a Lie subgroup of G (§ 1, no. 7, Proposition 14). The orbital mapping of \(\pi(e)\) is \(\pi \circ \phi\) . By Proposition 14 of § 1, L(G') is the kernel of L( \(\pi \circ \phi\) ) = \(T_e(\pi)\) \(\circ\) L( \(\phi\) ). The kernel of \(T_e(\pi)\) is L(H') (§ 1, no. 6, Proposition 11 (i)) and hence Ker L( \(\pi \circ \phi\) ) = L( \(\phi\) ) \(^{-1}\) (L(H')).

COROLLARY 1. Let \(\mathbf{G}\) , \(\mathbf{H}\) be Lie groups and \(\phi_1\) and \(\phi_2\) morphisms of \(\mathbf{G}\) into \(\mathbf{H}\) . Suppose that \(\mathbf{K}\) is of characteristic 0 and that \(\mathbf{H}\) is finite-dimensional. The set of \(g \in \mathbf{G}\) such that \(\phi_1(g) = \phi_2(g)\) is a Lie subgroup \(\mathbf{G}'\) of \(\mathbf{G}\) and \(\mathbf{L}(\mathbf{G}')\) is the set of \(x \in \mathbf{L}(\mathbf{G})\) such that \(\mathrm{L}(\phi_1)x = \mathrm{L}(\phi_2)x\) .

We write \(\phi(g) = (\phi_1(g), \phi_2(g))\) for all \(g \in G\) , so that \(\phi\) is a morphism of \(G\) into \(H \times H\) . Let \(\Delta\) be the diagonal subgroup of \(H \times H\) . Then \(G' = \phi^{-1}(\Delta)\) and \(L(\phi)x = (L(\phi_1)x, L(\phi_2)x)\) for all \(x \in L(G)\) . It now suffices to apply Proposition 29.

COROLLARY 2. Let \(\mathbf{G}\) be a finite-dimensional Lie group and \(\mathbf{G}_1\) and \(\mathbf{G}_2\) two Lie subgroups of \(\mathbf{G}\) . Suppose that \(\mathbf{K}\) is of characteristic 0. Then \(G_1 \cap G_2\) is a Lie subgroup of \(\mathbf{G}\) with Lie algebra \(\mathbf{L}(\mathbf{G}_1) \cap \mathbf{L}(\mathbf{G}_2)\) .

We apply Proposition 29 to the canonical injection of \(\mathbf{G}_1\) into \(\mathbf{G}\) and the subgroup \(\mathbf{G}_2\) .

COROLLARY 3. Let G, G', H be Lie groups and \(\phi: G \to H\) and \(\phi': G' \to H\) Lie group morphisms. Suppose that K is of characteristic 0 and that H is finite-dimensional. Let F be the set of \((g, g') \in G \times G'\) such that \(\phi(g) = \phi'(g')\) . Then F is a Lie subgroup of \(G \times G'\) and L(F) is the set of \((x, x') \in L(G) \times L(G')\) such that L( \(\phi\) ) \(x = L(\phi')x'\) .

We apply Corollary 1 to the morphisms \((g, g') \mapsto \phi(g)\) and \((g, g') \mapsto \phi'(g')\) of \(G \times G'\) into \(H\) .

PROPOSITION 30. Let G be a finite-dimensional Lie group with a countable base and

H and H' Lie subgroups of G. Suppose that K is of characteristic 0 and that HH' is locally closed in G.

\begin{enumerate}
\def\labelenumi{(\roman{enumi})}
\item
  HH' is a submanifold of G and \(\mathrm{T}_{e}(\mathrm{HH}') = \mathrm{L}(\mathrm{H}) + \mathrm{L}(\mathrm{H}')\) .
\item
  Suppose that every element of H commutes with every element of H'. Then HH' is a Lie subgroup of G. Let \(\phi\) be the mapping \((h, h') \mapsto hh'\) of H \(\times\) H' onto HH'. The kernel of \(\phi\) is the set of \((m, m^{-1})\) where \(m \in \mathrm{H} \cap \mathrm{H}'\) and the morphism of \((\mathrm{H} \times \mathrm{H}') / \mathrm{Ker} \phi\) onto HH' derived from \(\phi\) by passing to the quotient is a Lie group isomorphism.
\end{enumerate}

Let \(\mathbf{H} \times \mathbf{H}'\) operate on \(\mathbf{G}\) on the right by the mapping \(((h, h'), g) \mapsto hgh'^{-1}\) . The orbital mapping \(\rho\) of \(e\) is \((h, h') \mapsto hh'^{-1}\) . By Proposition 14 (iii) of § 1, no. 7, HH' is a submanifold of \(\mathbf{G}\) and \(\mathrm{T}_e(\mathrm{HH}') = \mathrm{Im} \mathrm{T}_e(\rho)\) . Now

\[
\mathrm{T} _ {e} (\rho) (\mathrm{L} (\mathrm{H}) \times \{0 \}) = \mathrm{L} (\mathrm{H}) \quad \text { and } \quad \mathrm{T} _ {e} (\rho) (\{0 \} \times \mathrm{L} (\mathrm{H} ^ {\prime})) = \mathrm{L} (\mathrm{H} ^ {\prime})
\]

and hence \(\mathrm{T}_{e}(\mathrm{HH}^{\prime})=\mathrm{L}(\mathrm{H})+\mathrm{L}(\mathrm{H}^{\prime})\) . Suppose that every element of H commutes with every element of \(H^{\prime}\) . Then \(HH^{\prime}\) is a subgroup of G. By (i), it is a Lie subgroup of G. The rest of the statement follows from Proposition 28.

PROPOSITION 31. Let G be a finite-dimensional Lie group with countable base, H a normal Lie subgroup of G and A a Lie subgroup of G. Suppose that K is of characteristic 0 and that AH is closed. Let \(\phi\) be the canonical morphism of G onto G/H. Then the canonical mappings

\[
\mathrm{A} / (\mathrm{H} \cap \mathrm{A}) \rightarrow \phi (\mathrm{A}), \quad \mathrm{AH} / \mathrm{H} \rightarrow \phi (\mathrm{A})
\]

are Lie group isomorphisms.

By Proposition 30, AH is a Lie subgroup of G. By Corollary 2 to Proposition 29, H ∩ A is a Lie subgroup of G. It is therefore meaningful to speak of the groups AH/H and A/(H ∩ A). On the other hand, φ(A), which is the canonical image of AH in G/H, is closed and hence is a Lie subgroup of G/H (Proposition 28 (iii)). Proposition 28, applied to the composite morphisms A → G → G/H and AH → G → G/H, proves that the canonical mappings of the proposition are Lie group isomorphisms.

PROPOSITION 32. Let G and H be Lie groups, k a non-discrete closed subfield of K and \(\phi\) a morphism of G into H for the Lie group structures over \(k\) . Suppose that K is of characteristic 0. If \(L(\phi)\) is K-linear, \(\phi\) is a morphism for the Lie group structures over K.

For all \(g \in G\) ,

\[
\mathrm{T} _ {g} (\phi) = \mathrm{T} _ {e} (\gamma (\phi (g))) \circ \mathrm{L} (\phi) \circ \mathrm{T} _ {g} (\gamma (g) ^ {- 1})
\]

and hence \(\mathrm{T}_g(\phi)\) is K-linear. The proposition then follows from Differentiable and Analytic Manifolds, R, 5.14.6.

